% ACL/EMNLP require an honest Limitations section.

\paragraph{Scope of evaluation.} The main evaluation is intentionally
narrow: LongMemEval, LoCoMo, and LongMemEval-V2 probe memory and
context selection, not general agent competence. We do not claim
coverage of code-generation agents, web/desktop control agents,
multi-lingual agents, or embodied agents. Generalisation to those
regimes remains to be shown, and in particular the
\textsc{glance}/\textsc{overview}/\textsc{detail} schema may need
domain-specific defaults for tasks where a one-line summary is not
informative (e.g., spatial reasoning, numerical tables).

\paragraph{Reliance on agent judgement.} The effectiveness of
multi-level disclosure depends on the agent choosing an appropriate
level per query. A poorly prompted agent that always requests
\textsc{detail} will obtain the token profile of the flat-window
baseline. We expose the choice explicitly rather than hiding it, on
the assumption that transparent control is preferable to a learned
but opaque policy; the assumption merits its own study.

\paragraph{Path-schema engineering.} Paths are conventions, not
contracts. A workspace with inconsistent path hygiene (for example,
similar content placed under \texttt{/notes/\dots} by one user and
\texttt{/memos/\dots} by another) reduces the benefit of prefix and
glob queries. We provide defaults but do not enforce them; larger
deployments may want path linting or migration tools that we do not
study here.

\paragraph{Infrastructure cost.} The audit log is append-only and
persistent. For very high-volume agents the PostgreSQL event table
grows quickly, and we have not evaluated long-term retention policies
(pruning, tiering to cold storage, compaction). Running
\textsc{Structure} at scale also assumes an operational Redis and
PostgreSQL, a heavier dependency footprint than a single-process
vector store.

\paragraph{LLM dependence.} All results are conditioned on the
specific backbone model used. We expect the relative ordering of
methods to transfer across models of similar scale, but smaller or
weaker models may lack the self-awareness to pick disclosure levels
well, narrowing the advantage over a simpler baseline.

\paragraph{Evaluation methodology.} Agent benchmarks are still
maturing. Task success is often graded heuristically, token counts
depend on tokenizer choice, and wall-clock latency reflects our
specific infrastructure. We report all three anyway, and release
configurations so that replications can recompute on their own
tokenizers and hardware. The hot-research LongMemEval-V2 result is a
fixture-level harness check rather than a full benchmark result; it
should not be interpreted as evidence of improved full-dataset
accuracy until the upstream split and evaluator are run.

\paragraph{Rating signal is agent-generated.} Source ratings are
produced by the same agent that consumes them, which introduces a
feedback loop: a confidently wrong rating is persisted and will bias
future retrieval until it is overwritten. We mitigate by event-sourcing
ratings (so offline review is possible) and by bounding $\alpha$ below
$1$ so similarity always remains a second signal, but we do not study
adversarial or self-reinforcing rating collapse, which is a known risk
of verbal-self-critique systems.

\paragraph{GC policy is hand-tuned.} The default event-TTL table is
derived from operational experience rather than learned from data.
Choosing TTLs that are too short drops tool results the agent still
needs; too long reintroduces the middle-of-context problem we aimed to
avoid. Per-type TTL is a coarser knob than per-event importance, and
an agent-learned GC policy is an obvious next step we have not
evaluated.
